\documentclass[10pt,a4paper]{amsart}
\usepackage[margin=19mm]{geometry}
\usepackage{amsmath,amssymb,mathtools}
\usepackage[scr=rsfs,cal=euler]{mathalfa}
\usepackage{xfrac}
\usepackage[unicode,hidelinks,bookmarks=false]{hyperref}
\newcommand{\ie}{i.e.\ }
\newcommand{\cf}{cf.\ }
\newcommand{\resp}{respectively\ }
\newcommand{\Subsection}{\subsection}
\providecommand{\nobreakdash}{\nobreak\hbox{-}}
\setlength{\emergencystretch}{2em}
\multlinegap0pt
\usepackage{bm}
\usepackage{bbm}
\usepackage{dsfont}
\usepackage{xspace}
\usepackage{cancel}
\usepackage{mathtools}
\usepackage{amsthm}
\makeatletter
\@ifundefined{theorem}{\newtheorem{theorem}{Theorem}}{}
\@ifundefined{corollary}{\newtheorem{corollary}[theorem]{Corollary}}{}
\@ifundefined{proposition}{\newtheorem{proposition}[theorem]{Proposition}}{}
\makeatother
\DeclareMathOperator{\Mag}{Mag}
\title[Small-Scale Magnitude Below One for Cyclic Two-Chunk Finite ]{Small-Scale Magnitude Below One for Cyclic Two-Chunk Finite Metric Spaces}
\author{Tsubasa Kamiyama}
\date{Source: arXiv:2606.00707v1 (CC BY 4.0)}
\begin{document}
\maketitle

\noindent\emph{Source and licence.} Small-Scale Magnitude Below One for Cyclic Two-Chunk Finite Metric Spaces. Version arXiv:2606.00707v1 (CC BY 4.0), distributed under the Creative Commons Attribution 4.0 International licence, as recorded in the arXiv licence metadata for this exact version and shown on the abstract page https://arxiv.org/abs/2606.00707v1. Full rights notice: licenses/arxiv-2606.00707-CC-BY-4.0.txt.\par\smallskip

\noindent\textit{Editorial scope.} Counted unit: the Proposition whose source label is \texttt{prop:n-3-4-5} in the article (source0.tex lines 590--595, source label \texttt{prop:n-3-4-5}), delivered together with the article's own complete original proof of it (source0.tex lines 597--839, one \texttt{proof} environment of 243 lines). No mathematics is written, repaired or re-derived: statement and proof are the article's own text line for line. Required context retained uncounted: eq:magnitude-formula (lines 375--377); cor:one-plus-formula (lines 537--550). Every definition, display and label of those lines is retained unchanged. Omitted from this package and not redistributed: 1--374, 378--536, 551--589, 596--596, 840--1096. Nothing inside a retained excerpt is omitted or altered: every retained body is delivered exactly as the article's lines are written, so no omission receipt of any kind accompanies this package. Citations: the retained excerpts carry no citation command at all, so no bibliography entry of the article is transcribed and no reference is redistributed. Reference adaptation: none was needed and none was made; every reference inside the delivered unit resolves inside it, so no unresolved reference or citation is produced. Printed slips kept literal: no printed word, symbol, hypothesis or number of the counted statement or of its proof was altered. Layout and class: the article's own class is \texttt{amsart} with packages including amsmath, amssymb, amsthm, mathtools, enumitem, graphicx, hyperref; the wrapper is the lane's pinned \texttt{amsart} editorial wrapper at 10pt on A4, which restates the theorem-like environments the retained text needs and adds the article's own macro definitions that the retained text needs (\texttt{\textbackslash Mag}), with extra wrapper packages (bm, bbm, dsfont, xspace, cancel, mathtools). The delivered numbering is the unit's own. Assets: no retained span contains a figure environment, a graphics-inclusion command, a PDF-page inclusion, a table or a TikZ picture; no image file is copied into this package, the package declares no asset dependency and no image file is redistributed with it. Licence: version arXiv:2606.00707v1 of this article is distributed under the Creative Commons Attribution 4.0 International licence, as recorded in the arXiv licence metadata for this exact version and shown on the abstract page https://arxiv.org/abs/2606.00707v1. A full rights notice is delivered at licenses/arxiv-2606.00707-CC-BY-4.0.txt. Characters: every retained excerpt is pure ASCII, so the delivered engine, XeLaTeX with its default Latin Modern text font, reports no missing character.
\begin{equation}\label{eq:magnitude-formula}
\Mag(tX)=\frac{n\bigl(A(t)+B(t)-2C(t)\bigr)}{A(t)B(t)-C(t)^2}.
\end{equation}

\begin{corollary}[Explicit balanced small-scale formula]\label{cor:one-plus-formula}
Let \(tX\) be a balanced \(t\)-scaled \((n,k;\alpha,\beta,\gamma,\delta)\)-two-chunk cyclic metric space.  Define
\begin{equation}\label{eq:Omega-definition}
\Omega_{n,k}(\alpha,\beta,\gamma,\delta)
=2(n-1)(\alpha^2+\beta^2)-4k(n-k)(\gamma-\delta)^2.
\end{equation}
If \(\Omega_{n,k}(\alpha,\beta,\gamma,\delta)\neq 0\), then
\begin{equation}\label{eq:one-plus-formula}
\lim_{t\to 0+}\Mag(tX)
=
1+\frac{(n-1)^2(\alpha-\beta)^2}
{\Omega_{n,k}(\alpha,\beta,\gamma,\delta)}.
\end{equation}
\end{corollary}

\begin{proposition}[The cases \(n=3,4,5\)]\label{prop:n-3-4-5}
Let \(3\leq n\leq 5\), and let \(1\leq k\leq \lfloor n/2\rfloor\).  For every balanced \(t\)-scaled \((n,k;\alpha,\beta,\gamma,\delta)\)-two-chunk cyclic metric space for which the small-scale limit is finite,
\begin{equation}\label{eq:n-3-4-5-geq-one}
\lim_{t\to 0+}\Mag(tX)\geq 1.
\end{equation}
\end{proposition}

\begin{proof}
The possible pairs \((n,k)\) with \(3\leq n\leq 5\) and
\(1\leq k\leq \lfloor n/2\rfloor\) are
\begin{equation}\label{eq:small-n-all-pairs}
(3,1),\qquad (4,1),\qquad (4,2),\qquad (5,1),\qquad (5,2).
\end{equation}
We check these cases one by one, using only the triangle inequalities and the
balanced condition.

Set
\begin{equation}\label{eq:u-v-eta}
u=\max\{\alpha,\beta\},
\qquad
v=\min\{\alpha,\beta\},
\qquad
\eta=|\gamma-\delta|.
\end{equation}
The mixed triangle inequalities give
\begin{equation}\label{eq:eta-leq-v}
\eta\leq v.
\end{equation}
Recall that
\begin{equation}\label{eq:Omega-small-proof}
\Omega_{n,k}
=
2(n-1)(u^2+v^2)-4k(n-k)\eta^2.
\end{equation}
By Corollary \ref{cor:one-plus-formula}, it is enough to show that
\(\Omega_{n,k}\) cannot be negative.  Indeed, if \(\Omega_{n,k}>0\), then
\[
\lim_{t\to 0+}\Mag(tX)
=
1+\frac{(n-1)^2(\alpha-\beta)^2}{\Omega_{n,k}}
\geq 1.
\]

We first treat the three cases with \(k=1\).

For \((n,k)=(3,1)\), we have
\begin{equation}\label{eq:n3k1-Omega}
\Omega_{3,1}=4(u^2+v^2)-8\eta^2.
\end{equation}
If \(\Omega_{3,1}<0\), then
\begin{equation}\label{eq:n3k1-contradiction}
\eta^2>\frac{u^2+v^2}{2}\geq v^2,
\end{equation}
hence \(\eta>v\), contradicting \eqref{eq:eta-leq-v}.  Thus
\(\Omega_{3,1}\geq 0\).

For \((n,k)=(4,1)\), we have
\begin{equation}\label{eq:n4k1-Omega}
\Omega_{4,1}=6(u^2+v^2)-12\eta^2.
\end{equation}
If \(\Omega_{4,1}<0\), then again
\begin{equation}\label{eq:n4k1-contradiction}
\eta^2>\frac{u^2+v^2}{2}\geq v^2,
\end{equation}
so \(\eta>v\), contradicting \eqref{eq:eta-leq-v}.  Hence
\(\Omega_{4,1}\geq 0\).

For \((n,k)=(5,1)\), we have
\begin{equation}\label{eq:n5k1-Omega}
\Omega_{5,1}=8(u^2+v^2)-16\eta^2.
\end{equation}
If \(\Omega_{5,1}<0\), then
\begin{equation}\label{eq:n5k1-contradiction}
\eta^2>\frac{u^2+v^2}{2}\geq v^2,
\end{equation}
so \(\eta>v\), contradicting \eqref{eq:eta-leq-v}.  Thus
\(\Omega_{5,1}\geq 0\).

Next consider \((n,k)=(4,2)\).  The balanced condition is
\begin{equation}\label{eq:n4k2-balanced-sum}
\gamma+\delta=\frac{3}{4}(u+v).
\end{equation}
Since \(k=2\) and \(n-k=2\), both the \(\gamma\)-edges and the
\(\delta\)-edges occur in pairs.  Therefore the triangle inequalities force
\begin{equation}\label{eq:n4k2-lower-bounds}
u\leq 2\gamma,
\qquad
u\leq 2\delta,
\end{equation}
or equivalently
\begin{equation}\label{eq:n4k2-gamma-delta-lower}
\gamma,\delta\geq \frac{u}{2}.
\end{equation}
In particular,
\begin{equation}\label{eq:n4k2-necessary}
\frac{3}{4}(u+v)=\gamma+\delta\geq u,
\end{equation}
so a metric space can exist only if \(u\leq 3v\).  Moreover, under the lower
bounds \eqref{eq:n4k2-gamma-delta-lower}, the largest possible value of
\(|\gamma-\delta|\) occurs when the smaller of \(\gamma,\delta\) is \(u/2\).
Thus
\begin{equation}\label{eq:n4k2-eta-bound}
\eta
\leq
\frac{3}{4}(u+v)-u
=
\frac{3v-u}{4}.
\end{equation}
Now
\begin{equation}\label{eq:n4k2-Omega}
\Omega_{4,2}=6(u^2+v^2)-16\eta^2.
\end{equation}
If \(\Omega_{4,2}<0\), then
\begin{equation}\label{eq:n4k2-negative-condition}
\eta^2>\frac{3}{8}(u^2+v^2).
\end{equation}
On the other hand, \eqref{eq:n4k2-eta-bound} gives
\begin{equation}\label{eq:n4k2-upper-square}
\eta^2\leq \frac{(3v-u)^2}{16}.
\end{equation}
But, since \(u\geq v>0\),
\begin{equation}\label{eq:n4k2-elementary-positive}
6(u^2+v^2)-(3v-u)^2
=
5u^2+6uv-3v^2
>0.
\end{equation}
Equivalently,
\begin{equation}\label{eq:n4k2-contradiction}
\frac{(3v-u)^2}{16}
<
\frac{3}{8}(u^2+v^2),
\end{equation}
which contradicts \eqref{eq:n4k2-negative-condition}.  Hence
\(\Omega_{4,2}>0\).

Finally consider \((n,k)=(5,2)\).  The balanced condition is
\begin{equation}\label{eq:n5k2-balanced}
2\gamma+3\delta=2(u+v).
\end{equation}
Again, since \(k=2\) and \(n-k=3\), the triangle inequalities force
\begin{equation}\label{eq:n5k2-lower-bounds}
u\leq 2\gamma,
\qquad
u\leq 2\delta,
\end{equation}
so
\begin{equation}\label{eq:n5k2-gamma-delta-lower}
\gamma,\delta\geq \frac{u}{2}.
\end{equation}
Therefore
\begin{equation}\label{eq:n5k2-necessary}
2(u+v)=2\gamma+3\delta\geq \frac{5u}{2},
\end{equation}
and hence a metric space can exist only if
\begin{equation}\label{eq:n5k2-u-bound}
u\leq 4v.
\end{equation}

We now bound \(\eta=|\gamma-\delta|\) directly.  If \(\gamma\geq \delta\),
then \(\gamma=\delta+\eta\), and \eqref{eq:n5k2-balanced} gives
\begin{equation}\label{eq:n5k2-case-one-delta}
5\delta+2\eta=2(u+v).
\end{equation}
Since \(\delta\geq u/2\), we obtain
\begin{equation}\label{eq:n5k2-case-one-eta}
\eta\leq v-\frac{u}{4}
=
\frac{4v-u}{4}.
\end{equation}
If \(\delta\geq \gamma\), then \(\delta=\gamma+\eta\), and
\eqref{eq:n5k2-balanced} gives
\begin{equation}\label{eq:n5k2-case-two-gamma}
5\gamma+3\eta=2(u+v).
\end{equation}
Since \(\gamma\geq u/2\), we obtain
\begin{equation}\label{eq:n5k2-case-two-eta}
\eta\leq \frac{4v-u}{6}
\leq
\frac{4v-u}{4},
\end{equation}
where the last inequality uses \eqref{eq:n5k2-u-bound}.  Hence, in all cases,
\begin{equation}\label{eq:n5k2-eta-bound}
\eta\leq \frac{4v-u}{4}.
\end{equation}

Now
\begin{equation}\label{eq:n5k2-Omega}
\Omega_{5,2}=8(u^2+v^2)-24\eta^2.
\end{equation}
If \(\Omega_{5,2}<0\), then
\begin{equation}\label{eq:n5k2-negative-condition}
\eta^2>\frac{u^2+v^2}{3}.
\end{equation}
But \eqref{eq:n5k2-eta-bound} implies
\begin{equation}\label{eq:n5k2-upper-square}
\eta^2\leq \frac{(4v-u)^2}{16}.
\end{equation}
Since \(u\geq v>0\),
\begin{equation}\label{eq:n5k2-elementary-positive}
16(u^2+v^2)-3(4v-u)^2
=
13u^2+24uv-32v^2
>0.
\end{equation}
Equivalently,
\begin{equation}\label{eq:n5k2-contradiction}
\frac{(4v-u)^2}{16}
<
\frac{u^2+v^2}{3},
\end{equation}
contradicting \eqref{eq:n5k2-negative-condition}.  Hence
\(\Omega_{5,2}>0\).

We have shown, case by case, that \(\Omega_{n,k}<0\) is impossible for every
pair in \eqref{eq:small-n-all-pairs}.  Therefore, in all cases with
\(\Omega_{n,k}>0\), Corollary \ref{cor:one-plus-formula} gives
\[
\lim_{t\to 0+}\Mag(tX)\geq 1.
\]

It remains only to handle the possible equality case \(\Omega_{n,k}=0\).
The above argument shows that equality can occur only in the cases \(k=1\),
and then necessarily
\begin{equation}\label{eq:k1-equality-case}
u=v,
\qquad
\eta=v.
\end{equation}
Thus \(\alpha=\beta=v\), and hence \(A(t)=B(t)\).  The magnitude formula
\eqref{eq:magnitude-formula} becomes, after cancelling the common factor
\(A(t)-C(t)\),
\begin{equation}\label{eq:k1-degenerate-cancellation}
\Mag(tX)
=
\frac{2n(A(t)-C(t))}{(A(t)-C(t))(A(t)+C(t))}
=
\frac{2n}{A(t)+C(t)}.
\end{equation}
Since \(A(0)=C(0)=n\), we get
\begin{equation}\label{eq:k1-degenerate-limit-final}
\lim_{t\to 0+}\Mag(tX)
=
\frac{2n}{2n}
=
1.
\end{equation}
Thus every balanced \(t\)-scaled \((n,k;\alpha,\beta,\gamma,\delta)\)-two-chunk cyclic metric space with
\(3\leq n\leq 5\) has finite small-scale magnitude at least \(1\).
\end{proof}

\end{document}
